\clearpage\fancyhead[L]{\small DETERMINATIO ATTRACTIONIS}\SourceLeaf{342}
{\small\begin{center}
DETERMINATIO ATTRACTIONIS\\
QUAM IN PUNCTUM QUODVIS POSITIONIS DATAE\\
EXERCERET PLANETA SI EIUS MASSA\\
PER TOTAM ORBITAM\\
RATIONE TEMPORIS QUO SINGULAE PARTES DESCRIBUNTUR\\
UNIFORMITER ESSET DISPERTITA.
\end{center}}

\begin{center}\rule{2cm}{0.3pt}\end{center}
\AttractionArticle{1}
Variationes saeculares, quas elementa orbitae planetariae a perturbatione
alius planetae patiuntur, ab huius positione in orbita sunt independentes, atque
eaedem forent, sive planeta perturbans in orbita elliptica secundum \textsc{Kepleri} leges
incedat, sive ipsius massa per orbitam eatenus aequabiliter dispertita concipiatur,
ut orbitae partibus, alias aequali temporis intervallo descriptis, iam aequales massae partes tribuantur, siquidem tempora revolutionum planetae perturbati et perturbantis non sint commensurabilia. Theorema hoc elegans, si a nemine hucusque disertis verbis propositum est, saltem perfacile ex astronomiae physicae principiis demonstratur. Problema itaque se offert tum per se, tum propter plura artificia, quae eius solutio requirit, attentione perdignum: attractionem orbitae planetariae, aut si mavis. annuli elliptici, cuius crassities infinite parva, atque secundum legem modo explicatam variabilis, in punctum quodlibet positione datum
exacte determinare.

\AttractionArticle{2}
Denotando excentricitatem orbitae per $e$, atque puncti cuiusvis in ipsa anomaliam excentricam per $E$, huius elemento $\mathrm{d}E$ respondebit elementum anomaliae mediae $(1-e\cos E)\mathrm{d}E$; quamobrem elementum massae ei orbitae portiunculae, cui respondent illa elementa, tribuendum, erit ad massam integram, quam
pro unitate accipiemus, ut $(1-e\cos E)\mathrm{d}E$ ad $2\pi$, exprimente $\pi$ semicircumfe\SourceLeaf{343}%
rentiam circuli pro radio 1. Statuendo itaque distantiam puncti attracti a puncto
orbitae $=\rho$, attractio ab orbitae elemento producta erit
\[
\begin{aligned}
&=\frac{(1-e\cos E)\mathrm{d}E}{2\pi\rho\rho}
\end{aligned}
\]
Designabimus semiaxem maiorem per $a$, semiaxem minorem per $b$, atque
illum tamquam lineam abscissarum, centrumque ellipsis tamquam initium adoptabimus. Hinc erit $aa-bb=aaee$, abscissa puncti orbitae $=a\cos E$, ordinata $=b\sin E$. Denique distantiam puncti attracti a plano orbitae denotabimus
per $C$, atque coordinatas reliquas axi maiori et minori parallelas per $A$ et $B$. His
ita praeparatis, attractio elementi orbitae decomponetur in duas axi maiori et minori parallelas atque tertiam plano orbitae normalem, puta
\[
\begin{aligned}
&\frac{(A-a\cos E)(1-e\cos E)\mathrm{d}E}{2\pi\rho^3}=\mathrm{d}\xi\\[6pt]
&\frac{(B-b\sin E)(1-e\cos E)\mathrm{d}E}{2\pi\rho^3}=\mathrm{d}\eta\\[6pt]
&\frac{C(1-e\cos E)\mathrm{d}E}{2\pi\rho^3}=\mathrm{d}\zeta
\end{aligned}
\]
ubi $\rho=\surd\bigl((A-a\cos E)^2+(B-b\sin E)^2+CC\bigr)$.

Integratis hisce differentialibus ab $E=0$ usque ad $E=360^{0}$, prodibunt
attractiones partiales $\xi,\eta,\zeta$ secundum directiones, directionibus coordinatarum
oppositas, e quibus attractio integra composita erit. et quas per methodum notam ad quaslibet alias directiones referre licebit.

\AttractionArticle{3}
Rei summa iam in eo versatur, ut introducta loco ipsius $E$ alia variabili,
quantitas radicalis in formam simpliciorem redigatur. Ad hunc finem statuemus
\[
\begin{aligned}
&\cos E=\frac{\alpha+\alpha'\cos T+\alpha''\sin T}{\gamma+\gamma'\cos T+\gamma''\sin T},\\[6pt]
&
\sin E=\frac{\beta+\beta'\cos T+\beta''\sin T}{\gamma+\gamma'\cos T+\gamma''\sin T}
\end{aligned}
\]
ubi autem novem coëfficientes $\alpha,\alpha',\alpha''$ etc. manifesto non sunt penitus arbitrarii,
sed certis conditionibus satisfacere debent, quas ante omnia perscrutari oportet.
Primo observamus, substitutionem eandem manere, si omnes coëfficientes per
eundem factorem multiplicentur, ita ut absque generalitatis detrimento uni ex ipsis
valorem determinatum tribuere, e. g. statuere liceret $\gamma=1$: attamen concinnitatis caussa omnes novem aliquantisper indefiniti maneant. Porro monemus, ex\SourceLeaf{344}%
cludi debere valores tales, ubi $\alpha,\alpha',\alpha''$ vel $\beta,\beta',\beta''$ ipsis $\gamma,\gamma'.\gamma''$ resp. proportionales essent: alioquin enim $E$ haud amplius indeterminata maneret. Nequeunt igitur $\gamma'\alpha''-\gamma''\alpha',\gamma''\alpha-\gamma\alpha'',\gamma\alpha'-\gamma'\alpha$ simul evanescere.

Manifesto coëfficientes $\alpha,\alpha',\alpha''$ etc. ita comparati esse debent, ut fiat indefinite
\[
\left.\begin{aligned}
 &(\alpha+\alpha'\cos T+\alpha''\sin T)^2\\[5pt]
+&(\beta+\beta'\cos T+\beta''\sin T)^2\\[5pt]
-&(\gamma+\gamma'\cos T+\gamma''\sin T)^2
\end{aligned}\right\}=0
\]
unde necessario haec functio habere debet formam
\[
\begin{aligned}
&k(\cos T^2+\sin T^2-1)
\end{aligned}
\]
Hinc colligimus sex aequationes conditionales
\[
\left.\begin{aligned}
-\alpha\alpha-\beta\beta+\gamma\gamma&=k\\[5pt]
-\alpha'\alpha'-\beta'\beta'+\gamma'\gamma'&=-k\\[5pt]
-\alpha''\alpha''-\beta''\beta''+\gamma''\gamma''&=-k\\[5pt]
-\alpha'\alpha''-\beta'\beta''+\gamma'\gamma''&=0\\[5pt]
-\alpha''\alpha-\beta''\beta+\gamma''\gamma&=0\\[5pt]
-\alpha\alpha'-\beta\beta'+\gamma\gamma'&=0
\end{aligned}\right\}\quad\mathrm{(I)}
\]
Ab his aequationibus pendent plures aliae, quas evolvere operae pretium
erit. Statuendo brevitatis caussa
\[
\begin{aligned}
&\alpha\beta'\gamma''+\alpha'\beta''\gamma+\alpha''\beta\gamma'
-\alpha\beta''\gamma'-\alpha'\beta\gamma''-\alpha''\beta'\gamma
\\[6pt]
&\qquad{}=\varepsilon\quad\mathinner{.\,.\,.\,.\,.\,.\,.\,.}\quad\text{(II)}
\end{aligned}
\]
e combinatione aequationum (I) facile derivantur novem sequentes:
\[
\left.\begin{aligned}
\varepsilon\alpha&=-k(\beta'\gamma''-\gamma'\beta'')\\[5pt]
\varepsilon\beta&=-k(\gamma'\alpha''-\alpha'\gamma'')\\[5pt]
\varepsilon\gamma&=+k(\alpha'\beta''-\beta'\alpha'')\\[5pt]
\varepsilon\alpha'&=+k(\beta''\gamma-\gamma''\beta)\\[5pt]
\varepsilon\beta'&=+k(\gamma''\alpha-\alpha''\gamma)\\[5pt]
\varepsilon\gamma'&=-k(\alpha''\beta-\beta''\alpha)\\[5pt]
\varepsilon\alpha''&=+k(\beta\gamma'-\gamma\beta')\\[5pt]
\varepsilon\beta''&=+k(\gamma\alpha'-\alpha\gamma')\\[5pt]
\varepsilon\gamma''&=-k(\alpha\beta'-\beta\alpha')
\end{aligned}\right\}\quad\mathrm{(III)}
\]
\SourceLeaf{345}
E tribus primis harum aequationum rursus deducimus hanc:
\[
\begin{aligned}
&\varepsilon\alpha(\beta'\gamma''-\gamma'\beta'')+\varepsilon\beta(\gamma'\alpha''-\alpha'\gamma'')+\varepsilon\gamma(\alpha'\beta''-\beta'\alpha'')\\[6pt]
&=-k(\beta'\gamma''-\gamma'\beta'')^2-k(\gamma'\alpha''-\alpha'\gamma'')^2\\[6pt]
&\qquad{}+k(\alpha'\beta''-\beta'\alpha'')^2
\end{aligned}
\]
cui aequivalens est haec:
\[
\begin{aligned}
&\varepsilon\varepsilon=k(-\alpha'\alpha'-\beta'\beta'+\gamma'\gamma')(-\alpha''\alpha''-\beta''\beta''+\gamma''\gamma'')\\[6pt]
&\qquad{}-k(-\alpha'\alpha''-\beta'\beta''+\gamma'\gamma'')^2
\end{aligned}
\]
quae adiumento aequationum 2, 3, 4 in (I) mutatur in hanc:
\[
\begin{aligned}
&\varepsilon\varepsilon=k^3\;\ldotp\;\ldotp\;\ldotp\;\ldotp\;\ldotp\;\ldotp\;\ldotp\;\ldotp\;\ldotp\;\ldotp\;\ldotp\;\text{(IV)}
\end{aligned}
\]
Aeque facile ex aequationibus (I) derivantur hae:
\[
\left.\begin{aligned}
(\beta'\gamma''-\gamma'\beta'')^2&=-k(k-\alpha'\alpha'-\alpha''\alpha'')\\[5pt]
(\gamma'\alpha''-\alpha'\gamma'')^2&=-k(k-\beta'\beta'-\beta''\beta'')\\[5pt]
(\alpha'\beta''-\beta'\alpha'')^2&=+k(k+\gamma'\gamma'+\gamma''\gamma'')\\[5pt]
(\beta''\gamma-\gamma''\beta)^2&=+k(k+\alpha\alpha-\alpha''\alpha'')\\[5pt]
(\gamma''\alpha-\alpha''\gamma)^2&=+k(k+\beta\beta-\beta''\beta'')\\[5pt]
(\alpha''\beta-\beta''\alpha)^2&=-k(k-\gamma\gamma+\gamma''\gamma'')\\[5pt]
(\beta\gamma'-\gamma\beta')^2&=+k(k+\alpha\alpha-\alpha'\alpha')\\[5pt]
(\gamma\alpha'-\alpha\gamma')^2&=+k(k+\beta\beta-\beta'\beta')\\[5pt]
(\alpha\beta'-\beta\alpha')^2&=-k(k-\gamma\gamma+\gamma'\gamma')
\end{aligned}\right\}\quad\text{(V)}
\]

Exempli caussa evolutionem primae adscribimus, ad cuius instar reliquae
facile formabuntur. Aequationes 4, 2, 3 in (I) scilicet suppeditant
\[
\begin{aligned}
&(\gamma'\gamma''-\beta'\beta'')^2-(\gamma'\gamma'-\beta'\beta')(\gamma''\gamma''-\beta''\beta'')\\[6pt]
&\qquad{}=\alpha'\alpha'\alpha''\alpha''-(\alpha'\alpha'-k)(\alpha''\alpha''-k)
\end{aligned}
\]
quae aequatio evoluta protinus ipsam primam in (V) sistit.

Ex his aequationibus (V) concludimus, valorem $k=0$ in disquisitione
nostra haud admissibilem esse; hinc enim omnes novem quantitates $\beta'\gamma''-\gamma'\beta''$ etc.
necessario evanescerent, i. e. coëfficientes $\alpha$, $\alpha'$, $\alpha''$ tum ipsis $\beta$, $\beta'$, $\beta''$, tum ipsis
$\gamma$, $\gamma'$, $\gamma''$ proportionales evaderent. Hinc etiam, propter aequationem IV, quantitas $\varepsilon$ evanescere nequit; quamobrem $k$ necessario debet esse quantitas positiva,
siquidem omnes coëfficientes $\alpha$, $\alpha'$, $\alpha''$ etc. debent esse reales. Combinatis tribus
aequationibus primis in (III) cum tribus primis in (V), hae novae prodeunt, quae
manifesto a valore ipsius $k$ non evanescente pendent:
\SourceLeaf{346}
\[
\left.\begin{aligned}
\alpha\alpha-\alpha'\alpha'-\alpha''\alpha''&=-k\\[5pt]
\beta\beta-\beta'\beta'-\beta''\beta''&=-k\\[5pt]
\gamma\gamma-\gamma'\gamma'-\gamma''\gamma''&=+k
\end{aligned}\right\}\quad\text{(VI)}
\]
Combinatio reliquarum easdem produceret. His denique adiungimus tres sequentes:
\[
\left.\begin{aligned}
\beta\gamma-\beta'\gamma'-\beta''\gamma''&=0\\[5pt]
\gamma\alpha-\gamma'\alpha'-\gamma''\alpha''&=0\\[5pt]
\alpha\beta-\alpha'\beta'-\alpha''\beta''&=0
\end{aligned}\right\}\quad\text{(VII)}
\]
quae facile ex aequationibus III derivantur; e. g. secunda, quinta et octava suppeditant:
\[
\begin{aligned}
&\varepsilon\beta\gamma-\varepsilon\beta'\gamma'-\varepsilon\beta''\gamma''=-k\gamma(\gamma'\alpha''-\alpha'\gamma'')\\[6pt]
&\qquad{}-k\gamma'(\gamma''\alpha-\alpha''\gamma)-k\gamma''(\gamma\alpha'-\alpha\gamma')=0
\end{aligned}
\]
Manifesto hae quoque aequationes ab exclusione valoris $k=0$ sunt dependentes\HistoricalNote{25}.

Quoniam, ut iam supra monuimus, omnes coëfficientes $\alpha$, $\alpha'$, $\alpha''$ etc. per
eundem factorem multiplicare licet, unde valor ipsius $k$ per quadratum eiusdem
factoris multiplicatus prodibit, abhinc semper supponemus
\[
\begin{aligned}
&k=1
\end{aligned}
\]
quo pacto necessario quoque erit vel $\varepsilon=+1$ vel $\varepsilon=-1$. Patet itaque, novem coëfficientes $\alpha$, $\alpha'$, $\alpha''$ etc., inter quos sex aequationes conditionales adsunt,
ad tres quantitates ab invicem independentes reducibiles esse debere, quod quidem commodissime per tres angulos sequenti modo efficitur:
\[
\begin{aligned}
&\alpha=\cos L\operatorname{tang}N\\[6pt]
&\beta=\sin L\operatorname{tang}N\\[6pt]
&\gamma=\sec N\\[6pt]
&\alpha'=\cos L\cos M\sec N\pm\sin L\sin M\\[6pt]
&\beta'=\sin L\cos M\sec N\mp\cos L\sin M\\[6pt]
&\gamma'=\cos M\operatorname{tang}N\\[6pt]
&\alpha''=\cos L\sin M\sec N\mp\sin L\cos M\\[6pt]
&\beta''=\sin L\sin M\sec N\pm\cos L\cos M\\[6pt]
&\gamma''=\sin M\operatorname{tang}N
\end{aligned}
\]

\SourceLeaf{347}\noindent ubi signorum ambiguorum superiora referuntur ad casum $\varepsilon=+1$, inferiora ad
casum $\varepsilon=-1$. Attamen tractatio analytica ad maximam partem elegantius
sine usu horum angulorum absolvitur. Ceterum haud difficile foret, significationem geometricam tum horum angulorum, tum reliquarum quantitatum auxiliarium in hac disquisitione occurrentium assignare; hanc vero interpretationem ad
institutum nostrum haud necessariam lectori perito explicandam linquimus.

\AttractionArticle{4}

Si iam in expressione distantiae $\rho$ pro $\cos E$ et $\sin E$ valores supra assumti
substituuntur, illa in hanc formam transibit:
\[
\rho=\frac{\surd\left(\begin{gathered}G+G'\cos T^2+G''\sin T^2\\[5pt]
+2H\cos T\sin T+2H'\sin T+2H''\cos T\end{gathered}\right)}{\gamma+\gamma'\cos T+\gamma''\sin T}
\]
ubi coëfficientes $\alpha$, $\alpha'$, $\alpha''$ etc. ita determinabimus, ut salvis sex aequationibus
conditionalibus
\[
\left.\begin{aligned}
-\alpha\alpha-\beta\beta+\gamma\gamma&=1\\[5pt]
-\alpha'\alpha'-\beta'\beta'+\gamma'\gamma'&=-1\\[5pt]
-\alpha''\alpha''-\beta''\beta''+\gamma''\gamma''&=-1\\[5pt]
-\alpha'\alpha''-\beta'\beta''+\gamma'\gamma''&=0\\[5pt]
-\alpha''\alpha-\beta''\beta+\gamma''\gamma&=0\\[5pt]
-\alpha\alpha'-\beta\beta'+\gamma\gamma'&=0
\end{aligned}\right\}\quad\text{[1]}
\]
adeoque etiam reliquis inde demanantibus, fiat
\[
\begin{aligned}
&H=0,\quad H'=0,\quad H''=0
\end{aligned}
\]
quo pacto problema generaliter loquendo erit determinatum. Quodsi itaque denominatorem ipsius $\rho$ per $t$ denotamus, transire debet functio trium quantitatum $t$, $t\cos E$, $t\sin E$ haec
\[
\begin{aligned}
&(AA+BB+CC)tt+aa(t\cos E)^2+bb(t\sin E)^2-2aAt\, .\, t\cos E\\[6pt]
&\qquad{}-2bBt\, .\, t\sin E
\end{aligned}
\]
per substitutionem
\[
\begin{aligned}
&t\cos E=\alpha+\alpha'\cos T+\alpha''\sin T\\[6pt]
&t\sin E=\beta+\beta'\cos T+\beta''\sin T\\[6pt]
&t=\gamma+\gamma'\cos T+\gamma''\sin T
\end{aligned}
\]
\SourceLeaf{348}
in
\[
\begin{aligned}
&G+G'\cos T^2+G''\sin T^2
\end{aligned}
\]
Manifesto hoc idem est, ac si dicas, functionem trium indeterminatarum $x,y,z$
hanc $(W)$
\[
\begin{aligned}
&aaxx+bbyy+(AA+BB+CC)zz-2aAxz-2bByz
\end{aligned}
\]
per substitutionem
\[
\begin{aligned}
&x=\alpha u+\alpha'u'+\alpha''u''\\[6pt]
&y=\beta u+\beta'u'+\beta''u''\\[6pt]
&z=\gamma u+\gamma'u'+\gamma''u''
\end{aligned}
\]
in functionem indeterminatarum $u,u',u''$ hanc
\[
\begin{aligned}
&Guu+G'u'u'+G''u''u''
\end{aligned}
\]
transire debere. At quum ex his formulis, adiumento aequationum [1], facile
sequatur
\[
\begin{aligned}
&u=-\alpha x-\beta y+\gamma z\\[6pt]
&u'=\alpha'x+\beta'y-\gamma'z\\[6pt]
&u''=\alpha''x+\beta''y-\gamma''z
\end{aligned}
\]
manifesto functio $W$ identica esse debebit cum hac
\[
\begin{aligned}
&G(-\alpha x-\beta y+\gamma z)^2
+G'(\alpha'x+\beta'y-\gamma'z)^2
\\[6pt]
&\qquad{}+G''(\alpha''x+\beta''y-\gamma''z)^2
\end{aligned}
\]
unde habemus sex aequationes
\[
\left.
\begin{aligned}
aa&=G\alpha\alpha+G'\alpha'\alpha'+G''\alpha''\alpha''\\[5pt]
bb&=G\beta\beta+G'\beta'\beta'+G''\beta''\beta''\\[5pt]
AA+BB+CC&=G\gamma\gamma+G'\gamma'\gamma'+G''\gamma''\gamma''\\[5pt]
bB&=G\beta\gamma+G'\beta'\gamma'+G''\beta''\gamma''\\[5pt]
aA&=G\gamma\alpha+G'\gamma'\alpha'+G''\gamma''\alpha''\\[5pt]
0&=G\alpha\beta+G'\alpha'\beta'+G''\alpha''\beta''
\end{aligned}
\right\}\quad [2]
\]

Ex his duodecim aequationibus [1] et [2] incognitas nostras $G,\allowbreak G',\allowbreak G'',\allowbreak \alpha,\allowbreak \alpha',\allowbreak \alpha''$
etc., determinare oportebit.
\SourceLeaf{349}
\AttractionArticle{5}

E combinatione aequationum [1] et [2] facile derivantur sequentes:
\[
\begin{aligned}
&-\alpha aa+\gamma aA=\alpha G\\[6pt]
&-\beta bb+\gamma bB=\beta G\\[6pt]
&\gamma(AA+BB+CC)-\alpha aA-\beta bB=\gamma G
\end{aligned}
\]
unde fit porro
\[
\begin{aligned}
&\alpha=\frac{\gamma aA}{aa+G}\quad\mathinner{.\,.\,.\,.\,.\,.\,.\,.\,.\,.\,.\,.\,.\,.\,.\,.\,.}\quad\text{[3]}\\[6pt]
&\beta=\frac{\gamma bB}{bb+G}\quad\mathinner{.\,.\,.\,.\,.\,.\,.\,.\,.\,.\quad.\,.\,.\,.\,.\,.}\quad\text{[4]}\\[6pt]
&AA+BB+CC-\frac{aaAA}{aa+G}-\frac{bbBB}{bb+G}=G
\end{aligned}
\]
Ultimam sic quoque exhibere possumus
\[
\begin{aligned}
&\frac{AA}{aa+G}+\frac{BB}{bb+G}+\frac{CC}{G}=1\quad\mathinner{.\,.\,.\,.\,.\,.}\quad\text{[5]}
\end{aligned}
\]

Perinde e combinatione aequationum [1] et [2] deducimus
\[
\begin{aligned}
&\alpha'aa-\gamma'aA=\alpha'G'\\[6pt]
&\beta'bb-\gamma'bB=\beta'G'\\[6pt]
&-\gamma'(AA+BB+CC)+\alpha'aA+\beta'bB=\gamma'G'
\end{aligned}
\]
atque hinc
\[
\begin{aligned}
&\alpha'{}=\frac{\gamma'aA}{aa-G'}\quad\mathinner{.\,.\,.\,.\,.\,.\,.\,.\,.\,.\,.\,.\,.\,.\,.\,,\,.\,.}\quad\text{[6]}\\[6pt]
&\beta'{}=\frac{\gamma'bB}{bb-G'}\quad\mathinner{.\,.\,.\,.\,.\,.\,.\,.\,.\,.\,.\,.\,.\,.\,.\,.\,.\,.}\quad\text{[7]}\\[6pt]
&\frac{AA}{aa-G'}+\frac{BB}{bb-G'}-\frac{CC}{G'}=1\quad\mathinner{.\,.\,.\,.\quad.\,.}\quad\text{[8]}
\end{aligned}
\]
et prorsus simili modo
\[
\begin{aligned}
&\alpha''{}=\frac{\gamma''aA}{aa-G''}\quad\mathinner{.\,.\,.\,.\,.\,.\,.\,.\,.\,.\,.\,.\,.\,.\,.\,.\,.\,.}\quad\text{[9]}\\[6pt]
&\beta''{}=\frac{\gamma''bB}{bb-G''}\quad\mathinner{.\,.\,.\,.\,.\,.\,.\,.\,.\,.\,.\,.\,.\,.\,.\,.\,.}\quad\text{[10]}\\[6pt]
&\frac{AA}{aa-G''}+\frac{BB}{bb-G''}-\frac{CC}{G''}=1\quad\mathinner{.\,.\,.\,.\,.\,.}\quad\text{[11]}
\end{aligned}
\]
Patet itaque, $G,-G',-G''$ esse radices aequationis
\[
\begin{aligned}
&\frac{AA}{aa+x}+\frac{BB}{bb+x}+\frac{CC}{x}=1\quad\mathinner{.\,.\,.\,.\,.\,.\,.\,.}\quad\text{[12]}
\end{aligned}
\]
\SourceLeaf{350}
quae rite evoluta ita se habet
\[
\begin{aligned}
&x^3-(AA+BB+CC-aa-bb)xx\\[6pt]
&\qquad{}+(aabb-aaBB-aaCC-bbAA-bbCC)x\\[6pt]
&-aabbCC=0\quad\mathinner{.\,.\,.\,.\,.\,.\,.\,.\,.\,.\,.\,.\,.\,.\,.\,.}\quad\text{[13]}
\end{aligned}
\]

\AttractionArticle{6}

Iam de indole huius aequationis cubicae sequentia sunt notanda.

I.\quad Ex aequationis termino ultimo \(-aabbCC\) concluditur, eam certe habere radicem unam realem, et quidem vel positivam, vel, si \(C=0\), cifrae aequalem. Denotemus hanc radicem realem non negativam per \(g\).

II.\quad Subtrahendo ab aequatione 12, ita exhibita
\[
\begin{aligned}
&x=\frac{AAx}{aa+x}+\frac{BBx}{bb+x}+CC
\end{aligned}
\]
hanc\textperiodcentered
\[
\begin{aligned}
&g=\frac{AAg}{aa+g}+\frac{BBg}{bb+g}+CC
\end{aligned}
\]
et dividendo per \(x-g\), oritur nova, duas reliquas radices complectens
\[
\begin{aligned}
&1=\frac{aaAA}{(aa+x)(aa+g)}+\frac{bbBB}{(bb+x)(bb+g)}
\end{aligned}
\]
quae rite ordinata et soluta suppeditat \([14]\)
\[
\begin{aligned}
&2x=\frac{aaAA}{aa+g}+\frac{bbBB}{bb+g}-aa-bb\\[6pt]
&\qquad{}\pm\surd\!\left(\begin{gathered}\left(aa-bb-\frac{aaAA}{aa+g}+\frac{bbBB}{bb+g}\right)^2
\\[6pt]+\frac{4aabbAABB}{(aa+g)(bb+g)}\end{gathered}\right)
\end{aligned}
\]
Haec expressio, quum quantitas sub signo radicali natura sua sit positiva, vel
saltem non negativa, monstrat, etiam duas reliquas radices semper fieri reales.

III.\quad Subtrahendo autem ab invicem aequationes istas sic exhibitas
\[
\begin{aligned}
&gx=\frac{AAgx}{aa+x}+\frac{BBgx}{bb+x}+gCC\\[6pt]
&gx=\frac{AAgx}{aa+g}+\frac{BBgx}{bb+g}+xCC
\end{aligned}
\]
et dividendo per \(g-x\), prodit aequatio duas reliquas radices continens in hacce
forma:
\[
\begin{aligned}
&0=\frac{AAgx}{(aa+g)(aa+x)}+\frac{BBgx}{(bb+g)(bb+x)}+CC
\end{aligned}
\]
cui manifesto, si \(g\) est quantitas positiva, per valorem positivum ipsius \(x\) satisfieri nequit. Unde concludimus, aequationem nostram cubicam radices positivas
plures quam unam habere non posse.
\SourceLeaf{351}\par
IV.\quad Quoties itaque \(0\) non est inter radices aequationis nostrae, aderunt
necessario radix una positiva cum duabus negativis. Quoties vero \(C=0\), adeoque \(0\) una radicum, reliquas complectetur aequatio
\[
\begin{aligned}
&xx-(AA+BB-aa-bb)x+aabb-aaBB-bbAA=0
\end{aligned}
\]
unde hae radices exprimentur per
\[
\begin{aligned}
&\frac12(AA+BB-aa-bb)\pm\frac12\surd\!\left((AA-BB-aa+bb)^2+4AABB\right)
\end{aligned}
\]
Tres casus hic iterum distinguere oportebit.

\textit{Primo} si terminus ultimus \(aabb-aaBB-bbAA\) est positivus (i.\ e.\ si
punctum attractum in plano ellipsis attrahentis \textit{intra} curvam iacet), ambae radices, quum reales esse debeant, eodem signo affectae erunt, adeoque quum simul
positivae esse nequeant, necessario erunt negativae. Ceterum hoc etiam independenter ab iis, quae iam demonstrata sunt, inde concludi potest, quod co\"efficiens medius, quem ita exhibere licet
\[
\begin{aligned}
&(aabb-aaBB-bbAA)\left(\frac1{aa}+\frac1{bb}\right)+\frac{bbAA}{aa}+\frac{aaBB}{bb}
\end{aligned}
\]
manifesto in hoc casu sit positivus.

\textit{Secundo}, si terminus ultimus est negativus, sive punctum attractum in plano
ellipsis \textit{extra} curvam situm, necessario altera radix positiva erit, altera negativa.

\textit{Tertio} autem, si terminus ultimus ipse evanesceret, sive punctum attractum
in ipsa ellipsis circumferentia iaceret, etiam radix secunda fieret \(=0\), atque tertia
\[
\begin{aligned}
&=-\frac{bbAA}{aa}-\frac{aaBB}{bb}
\end{aligned}
\]
i.\ e.\ negativa. Ceterum hunc casum, physice impossibilem, et in quo attractio
ipsa infinite magna evaderet, a disquisitione nostra, hocce saltem loco, excludemus.

\AttractionArticle{7}

Ad determinandos co\"efficientes \(\gamma,\gamma',\gamma''\), ex aequationibus \(1,3,4,6,7,9,10\)
invenimus 
